% Einheit 1 von 4 – Prozentsatz berechnen (Katalog Einheit 2)
\einheitenkopf[e1][1 Prozentsatz]{Einheit 1 von 4 · Prozentsatz berechnen}
\zweigzeile{Hier lernst du auszurechnen, wie viel Prozent ein Teil vom Ganzen ist · neu in diesem Jahr · P10 oft · baut auf: Prozent am Streifen (Nr.~1–3), Bruch auf den Nenner 100 (Nr.~4), Runden (Nr.~6), Dreisatz (Nr.~8)}
\setcounter{aufgabe}{11}

\begin{aufgabe}{Ich erkenne, was das Ganze ist.}
Unterstreiche das Ganze.
\begin{teile}
\teil Von 20 Kindern der Klasse fahren 8 mit dem Rad.
\teil Tom hat 12 von 15 Würfen getroffen.
\teil Im Glas sind 30 Bonbons. 9 davon sind rot.
\teil 45 Plätze sind besetzt. Das Kino hat 60 Plätze.
\teil Lea hat schon 180 Seiten gelesen. Das Buch hat 240 Seiten.
\end{teile}
\end{aufgabe}

\begin{aufgabe}{Ich kann einen Streifen in gleich große Teile einteilen.}
Teile den Streifen ein.

\swa{in $50$-\%-Schritte}{\streifenleer[0]}
\swa{in $25$-\%-Schritte}{\streifenleer[0]}
\swa{in $10$-\%-Schritte}{\streifenleer[0]}
\swa{in $20$-\%-Schritte}{\streifenleer[0]}
\end{aufgabe}

\verfahren{Anteil in Prozent am Streifen}
\begin{aufgabe}{Ich kann am Streifen ablesen, wie viel Prozent ein Teil vom Ganzen ist.}
Lies am Streifen ab, wie viel Prozent das sind.

\swa{3 von 10 Kindern: \leerfeld[\%]}{\streifen{30}{$0$}{10 Kinder}}
\swa{$25\,€$ von $50\,€$: \leerfeld[\%]}{\streifen{50}{$0\,€$}{$50\,€$}}
\swa{$10\,$kg von $40\,$kg: \leerfeld[\%]}{\streifen[4]{25}{$0\,$kg}{$40\,$kg}}
\swa{$4\,$m von $20\,$m: \leerfeld[\%]}{\streifen[5]{20}{$0\,$m}{$20\,$m}}
\end{aufgabe}

\verfahren{Anteil in Prozent durch Erweitern}
\begin{aufgabe}{Ich kann durch Erweitern ausrechnen, wie viel Prozent ein Teil vom Ganzen ist.}
Schreibe als Bruch mit dem Nenner 100 und dann in Prozent. Den Anteil in Prozent nennt man Prozentsatz.
\rechenplatz{2}

\swz{5 von 100 Losen gewinnen.}{\streifen[20]{5}{$0$}{100 Lose}}
\swz{15 von 100 Losen gewinnen.}{\streifen[20]{15}{$0$}{100 Lose}}
\swz{35 von 100 Losen gewinnen.}{\streifen[20]{35}{$0$}{100 Lose}}
\swz{60 von 100 Losen gewinnen.}{\streifen[20]{60}{$0$}{100 Lose}}
\swfrage{Was bleibt gleich, was ändert sich?}
\end{aufgabe}

\begin{aufgabe}{Ich kann durch Erweitern ausrechnen, wie viel Prozent ein Teil vom Ganzen ist – weiter.}
Schreibe als Bruch mit dem Nenner 100 und dann in Prozent.

\swz{13 von 50 Schülern}{$\dfrac{13}{50} = \dfrac{\hspace{8mm}}{100} = \rule{10mm}{0.4pt}\,\%$}
\swz{6 von 25 Schülern}{$\dfrac{6}{25} = \dfrac{\hspace{8mm}}{100} = \rule{10mm}{0.4pt}\,\%$}
\swz{7 von 20 Schülern}{$\dfrac{7}{20} = \dfrac{\hspace{8mm}}{100} = \rule{10mm}{0.4pt}\,\%$}
\swz[3]{12 von 40 Schülern}{$\dfrac{12}{40} = \dfrac{\hspace{8mm}}{10} = \dfrac{\hspace{8mm}}{100} = \rule{10mm}{0.4pt}\,\%$}
\end{aufgabe}

\verfahren{Anteil in Prozent mit dem Taschenrechner}
\begin{aufgabe}{Ich kann mit dem Taschenrechner ausrechnen, wie viel Prozent ein Teil vom Ganzen ist.}
Rechne mit dem Taschenrechner. Gib das Ergebnis in Prozent an.

\swz{27 von 60 Heften}{$\dfrac{27}{60} = \rule{12mm}{0.4pt} = \rule{10mm}{0.4pt}\,\%$}
\swz{$4{,}5\,$kg von $18\,$kg}{$\dfrac{4{,}5}{18} = \rule{12mm}{0.4pt} = \rule{10mm}{0.4pt}\,\%$}
\anweisung{Runde auf eine Stelle nach dem Komma.}
\swz{7 von 45 Heften}{$\dfrac{7}{45} = \rule{12mm}{0.4pt} \approx \rule{10mm}{0.4pt}\,\%$}
\swz{35 von 270 Heften}{$\dfrac{35}{270} = \rule{12mm}{0.4pt} \approx \rule{10mm}{0.4pt}\,\%$}
\end{aufgabe}

\verfahren{Anteil in Prozent aus einem Sachtext}
\begin{aufgabe}{Ich kann in einem Sachtext ausrechnen, wie viel Prozent ein Teil vom Ganzen ist.}
Berechne, wie viel Prozent es sind.

\swz[3]{Von 25 Kindern haben 9 ein Haustier. Wie viel Prozent sind das?}{\streifen[0]{0}{$0$}{\rule{12mm}{0.4pt}}}
\swz[3]{Im Chor singen 12 Mädchen und 8 Jungen. Wie viel Prozent der Sänger sind Jungen?}{\streifen[0]{0}{$0$}{\rule{12mm}{0.4pt}}}
\end{aufgabe}

\begin{aufgabe}{Ich kann in einem Sachtext ausrechnen, wie viel Prozent ein Teil vom Ganzen ist – weiter.}
Berechne, wie viel Prozent es sind.

\swz[3]{Eine Jacke kostete $80\,€$, jetzt kostet sie $60\,€$. Wie viel Prozent Rabatt (Preisnachlass) gibt es?}{\streifen[0]{0}{$0\,€$}{\rule{12mm}{0.4pt}}}
\swz[4]{Lena backt Waffeln. 3 Waffeln sind verbrannt, 19 sind gelungen. Wie viel Prozent aller Waffeln sind verbrannt? Runde auf eine Stelle nach dem Komma. (P10 2018 OS)}{}
\end{aufgabe}

\begin{aufgabe}{Ich kann zu einem Prozentsatz passende Zahlen finden.}
Gib einen Teil und ein Ganzes an, die passen.

\swa{$50\,\%$ sind \rule{10mm}{0.4pt} von \rule{10mm}{0.4pt}}{\streifen{50}{$0\,\%$}{$100\,\%$}}
\swa{$25\,\%$ sind \rule{10mm}{0.4pt} von \rule{10mm}{0.4pt}}{\streifen[4]{25}{$0\,\%$}{$100\,\%$}}
\swa{$10\,\%$ sind \rule{10mm}{0.4pt} von \rule{10mm}{0.4pt}}{\streifen{10}{$0\,\%$}{$100\,\%$}}
\swa{$75\,\%$ sind \rule{10mm}{0.4pt} von \rule{10mm}{0.4pt}}{\streifen[4]{75}{$0\,\%$}{$100\,\%$}}
\end{aufgabe}

\begin{aufgabe}{Ich finde den Fehler, wenn jemand ausrechnet, wie viel Prozent ein Teil ist.}
Finde den Fehler, benenne ihn und rechne richtig.

\swz[3]{12 von 48 Äpfeln sind faul. Wie viel Prozent sind faul? Ben rechnet so:}{\rechnung{48 : 12 &= 4 \\ \text{Anteil} &= 4\,\%}}
\swz[3]{Im Beet stehen 6 rote und 14 gelbe Tulpen. Wie viel Prozent sind rot? Lea rechnet so:}{\rechnung{6 : 14 &\approx 0{,}43 \\ \text{Anteil} &\approx 43\,\%}}
\end{aufgabe}

\begin{aufgabe}{Ich kann begründen, warum man durch das Ganze teilt.}
Begründe ohne zu rechnen.

\swfrage{Tim rechnet bei „3 von 4“ die Aufgabe $4 : 3$. Hat er recht? Begründe.}
\swfrage{Warum sind 5 von 20 mehr Prozent als 5 von 25? Begründe.}
\end{aufgabe}

\begin{aufgabe}{Ich kann mit Prozent vergleichen, wer besser trifft.}
Mia trifft beim Basketball 18 von 24 Freiwürfen. Sara trifft 21 von 30 Freiwürfen. Rechne und entscheide.

\swz{Wie viel Prozent trifft Mia?}{\streifen[0]{0}{$0$}{24 Würfe}}
\swz{Wie viel Prozent trifft Sara?}{\streifen[0]{0}{$0$}{30 Würfe}}
\swfrage{Sara hat mehr Treffer. Wer wirft trotzdem besser? Begründe.}
\end{aufgabe}

\uebersichtskasten{Prozentsatz: Teil geteilt durch Ganzes – das ist der Anteil, dann in Prozent schreiben. \\ 14 von 40: \quad $14 : 40 = 0{,}35 = 35\,\%$ \\ 3 von 8: \quad $3 : 8 = 0{,}375 = 37{,}5\,\%$}
